\documentclass{article}
\usepackage{amsmath}
\usepackage{amssymb}

\title{Gaussian Quadrature for Periodic Functions}
\author{}
\date{}

\begin{document}

\maketitle

Given a complex-valued function \(f\) satisfying \(f(-L) = f(L)\), we wish to compute,
\begin{equation}
    I = \int_{-L}^{L} f(t)  dt
\end{equation}

using Gaussian quadrature with \(M+1\) points:
\begin{equation}
I(M) = \sum_{j=0}^{M} f(\tau_j) w_j
\end{equation}

where \(\tau_j\) are nodes and \(w_j\) are weights.

We can do a coordinate transformation, by appling the substitution \(t = Lx\):
\begin{equation}
    I = \int_{-1}^{1} f(Lx) \cdot L  dx = L \int_{-1}^{1} g(x)  dx, \quad g(x) \equiv f(Lx)
\end{equation}

The periodic condition becomes \(g(-1) = g(1)\).

For integrals over \([-1, 1]\), use the \textbf{Gauss-Legendre quadrature} with \(n = M+1\) points. 

The nodes \(\tau_j\) are roots of the Legendre polynomial \(P_n(x)\) of degree \(n\), i.e.,
\begin{equation}
    P_n(\tau_j) = 0, \quad \forall j\in[0,..,M]
\end{equation}

and weights \(w_j\) are given by,
\begin{equation}
    w_j = \frac{2}{(1 - \tau_j^2) [P'_n(\tau_j)]^2} = \frac{2(1 - \tau_j^2)}{n^2 [P_{n-1}(\tau_j)]^2}
\end{equation}

The integral is then approximated by:
\[
I(M) = L \sum_{j=0}^{M} f(L \tau_j) w_j
\]
where \(\{\tau_j, w_j\}_{j=0}^{M}\) are the Gauss-Legendre nodes and weights for \([-1, 1]\) with \(n = M+1\) points.


\subsection*{Properties}
\begin{itemize}
  \item Symmetry: \(\tau_j = -\tau_{M-j}\),  \(w_j = w_{M-j}\)
  \item Exact for polynomials up to degree \(2n-1 = 2M+1\)
  \item Weights satisfy \(\sum_{j=0}^{M} w_j = 2\)
\end{itemize}


\end{document}